% GENERATED FILE. Edit canonical lessons, then run npm run generate:books.
% canonical-source-hash: fnv1a64:f7073c86482f0484
% canonical-lessons: JA-C292-oiru, JA-C292-toshiyori, JA-C292-jinsei, JA-C292-shinu, JA-C292-sasaeru

\chapter{A Long Life}
\label{ch:ja-a-long-life}
\hlchaptermodality{292}

\begin{chapteropening}
\textbf{By the end of this chapter:} \emph{I can say how people age and who looks after them.}

Complete the last lesson of chapter 292: I can say how people age and who looks after them.
\end{chapteropening}

\section[\texorpdfstring{\ja{おいる} (oiru)}{おいる (oiru)}]{\ja{おいる} (oiru) — to grow old; to age}

\label{lesson:JA-C292-oiru}



\begin{quote}
Before the new one: say the Japanese for a grave, then the Japanese for to pray.
\end{quote}

\subsection*{You'll want to know: \ja{おいる}}
\textbf{\ja{おいる}} — \emph{oiru} — \textquotedblleft{}to grow old; to age\textquotedblright{}.

\begin{quote}
\textbf{\ja{ちちも} \ja{すこし} \ja{おいました。}} — \emph{chichi mo sukoshi oimashita} — \textquotedblleft{}My father has aged a little too.\textquotedblright{}
\end{quote}

Polite \textbf{\ja{おいます}}, \emph{oimasu}; an ichidan verb. Past (polite): \textbf{\ja{そふが} \ja{おいました。}}, \emph{sofu ga oimashita}.

\begin{cousinweb}
\textbf{\ja{おいる}} (\emph{oiru}) is a native Japanese word (wago), inherited rather than borrowed from Chinese; its kanji was attached later as a kun'yomi (JMdict).
\end{cousinweb}

\begin{grammarlens}[title={What you've built}]
One more word for telling what happened.
\end{grammarlens}

\subsection*{Your turn}
\begin{itemize}
\raggedright
  \item \emph{Say it:} \emph{oiru}
  \item \emph{Say it:} \emph{oiru}, once more
  \item \emph{Say it:} \emph{chichi mo sukoshi oimashita} — My father has aged a little too.
  \item \emph{Recall:} say the Japanese for a grave, then the Japanese for to pray, then say \emph{oiru} again
\end{itemize}

\begin{tcolorbox}[breakable,colback=teal!4,colframe=teal!35!black,title={Before you move on}]
What does \ja{おいる} mean? (\textquotedblleft{}To grow old; to age\textquotedblright{}.) Say it once more.
\end{tcolorbox}
\section[\texorpdfstring{\ja{としより} (toshiyori)}{としより (toshiyori)}]{\ja{としより} (toshiyori) — an old person; the elderly}

\label{lesson:JA-C292-toshiyori}



\begin{quote}
Before the new one: say the Japanese for to pray, then the Japanese for to grow old.
\end{quote}

\subsection*{You'll want to know: \ja{としより}}
\textbf{\ja{としより}} — \emph{toshiyori} — \textquotedblleft{}an old person; the elderly\textquotedblright{}.

\begin{quote}
\textbf{\ja{その} \ja{まちには} \ja{としよりが} \ja{おおいです。}} — \emph{sono machi ni wa toshiyori ga ōi desu} — \textquotedblleft{}There are many old people in that town.\textquotedblright{}
\end{quote}

\begin{cousinweb}
\textbf{\ja{としより}} (\emph{toshiyori}) is a native Japanese word (wago), inherited rather than borrowed from Chinese, built from \emph{toshi}, \textquotedblleft{}years\textquotedblright{}, and \emph{yori}, \textquotedblleft{}gathering\textquotedblright{} (from \emph{yoru}) (JMdict, KANJIDIC2).
\end{cousinweb}

\begin{grammarlens}[title={What you've built}]
One more word for telling what happened.
\end{grammarlens}

\subsection*{Your turn}
\begin{itemize}
\raggedright
  \item \emph{Say it:} \emph{toshiyori}
  \item \emph{Say it:} \emph{toshiyori}, once more
  \item \emph{Say it:} \emph{sono machi ni wa toshiyori ga ōi desu} — There are many old people in that town.
  \item \emph{Recall:} say the Japanese for to pray, then the Japanese for to grow old, then say \emph{toshiyori} again
\end{itemize}

\begin{tcolorbox}[breakable,colback=teal!4,colframe=teal!35!black,title={Before you move on}]
What does \ja{としより} mean? (\textquotedblleft{}An old person; the elderly\textquotedblright{}.) Say it once more.
\end{tcolorbox}
\section[\texorpdfstring{\ja{じんせい} (jinsei)}{じんせい (jinsei)}]{\ja{じんせい} (jinsei) — a life; one's lifetime}

\label{lesson:JA-C292-jinsei}



\begin{quote}
Before the new one: say the Japanese for to grow old, then the Japanese for an old person.
\end{quote}

\subsection*{You'll want to know: \ja{じんせい}}
\textbf{\ja{じんせい}} — \emph{jinsei} — \textquotedblleft{}a life; one's lifetime\textquotedblright{}.

\begin{quote}
\textbf{\ja{日本人}\ja{の} \ja{じんせいは} \ja{ながく} \ja{なりました。}} — \emph{nihonjin no jinsei wa nagaku narimashita} — \textquotedblleft{}Japanese people's lives have grown longer.\textquotedblright{}
\end{quote}

\begin{cousinweb}
\textbf{\ja{じんせい}} (\emph{jinsei}) is Sino-Japanese (kango): two Chinese characters read with their on'yomi, readings borrowed from Chinese. They are \emph{jin}, \textquotedblleft{}person\textquotedblright{} (Middle Chinese \emph{nyin}) and \emph{sei}, \textquotedblleft{}life\textquotedblright{} (\emph{srjaeng}) (KANJIDIC2; Middle Chinese in Baxter's transcription).
\end{cousinweb}

\begin{grammarlens}[title={What you've built}]
One more word for telling what happened.
\end{grammarlens}

\subsection*{Your turn}
\begin{itemize}
\raggedright
  \item \emph{Say it:} \emph{jinsei}
  \item \emph{Say it:} \emph{jinsei}, once more
  \item \emph{Say it:} \emph{nihonjin no jinsei wa nagaku narimashita} — Japanese people's lives have grown longer.
  \item \emph{Recall:} say the Japanese for to grow old, then the Japanese for an old person, then say \emph{jinsei} again
\end{itemize}

\begin{tcolorbox}[breakable,colback=teal!4,colframe=teal!35!black,title={Before you move on}]
What does \ja{じんせい} mean? (\textquotedblleft{}A life; one's lifetime\textquotedblright{}.) Say it once more.
\end{tcolorbox}
\section[\texorpdfstring{\ja{しぬ} (shinu)}{しぬ (shinu)}]{\ja{しぬ} (shinu) — to die}

\label{lesson:JA-C292-shinu}



\begin{quote}
Before the new one: say the Japanese for an old person, then the Japanese for a life.
\end{quote}

\subsection*{You'll want to know: \ja{しぬ}}
\textbf{\ja{しぬ}} — \emph{shinu} — \textquotedblleft{}to die\textquotedblright{}.

\begin{quote}
\textbf{\ja{そふの} \ja{いぬが} \ja{きょねん} \ja{しにました。}} — \emph{sofu no inu ga kyonen shinimashita} — \textquotedblleft{}Grandfather's dog died last year.\textquotedblright{}
\end{quote}

Polite \textbf{\ja{しにます}}, \emph{shinimasu}; a godan verb. Past (polite): \textbf{\ja{とりが} \ja{しにました。}}, \emph{tori ga shinimashita}.

\begin{cousinweb}
\textbf{\ja{しぬ}} (\emph{shinu}) is a native Japanese word (wago), inherited rather than borrowed from Chinese; its kanji was attached later as a kun'yomi (JMdict).
\end{cousinweb}

\begin{grammarlens}[title={What you've built}]
One more word for telling what happened.
\end{grammarlens}

\subsection*{Your turn}
\begin{itemize}
\raggedright
  \item \emph{Say it:} \emph{shinu}
  \item \emph{Say it:} \emph{shinu}, once more
  \item \emph{Say it:} \emph{sofu no inu ga kyonen shinimashita} — Grandfather's dog died last year.
  \item \emph{Recall:} say the Japanese for an old person, then the Japanese for a life, then say \emph{shinu} again
\end{itemize}

\begin{tcolorbox}[breakable,colback=teal!4,colframe=teal!35!black,title={Before you move on}]
What does \ja{しぬ} mean? (\textquotedblleft{}To die\textquotedblright{}.) Say it once more.
\end{tcolorbox}
\section[\texorpdfstring{\ja{ささえる} (sasaeru)}{ささえる (sasaeru)}]{\ja{ささえる} (sasaeru) — to support; to hold up}

\label{lesson:JA-C292-sasaeru}



\begin{quote}
Before the new one: say the Japanese for a life, then the Japanese for to die.
\end{quote}

\subsection*{You'll want to know: \ja{ささえる}}
\textbf{\ja{ささえる}} — \emph{sasaeru} — \textquotedblleft{}to support; to hold up\textquotedblright{}.

\begin{quote}
\textbf{\ja{かぞくが} \ja{いつも} \ja{ちちを} \ja{ささえています。}} — \emph{kazoku ga itsumo chichi o sasaete imasu} — \textquotedblleft{}The family is always there for my father.\textquotedblright{}
\end{quote}

Polite \textbf{\ja{ささえます}}, \emph{sasaemasu}; an ichidan verb. Past (polite): \textbf{\ja{かぞくが} \ja{そぼを} \ja{ささえました。}}, \emph{kazoku ga sobo o sasaemashita}.

\begin{cousinweb}
\textbf{\ja{ささえる}} (\emph{sasaeru}) is a native Japanese word (wago), inherited rather than borrowed from Chinese; its kanji was attached later as a kun'yomi (JMdict).
\end{cousinweb}

\begin{grammarlens}[title={What you've built}]
One more word for telling what happened.
\end{grammarlens}

\subsection*{Your turn}
\begin{itemize}
\raggedright
  \item \emph{Say it:} \emph{sasaeru}
  \item \emph{Say it:} \emph{sasaeru}, once more
  \item \emph{Say it:} \emph{kazoku ga itsumo chichi o sasaete imasu} — The family is always there for my father.
  \item \emph{Recall:} say the Japanese for a life, then the Japanese for to die, then say \emph{sasaeru} again
\end{itemize}

\begin{tcolorbox}[breakable,colback=teal!4,colframe=teal!35!black,title={Before you move on}]
What does \ja{ささえる} mean? (\textquotedblleft{}To support; to hold up\textquotedblright{}.) Say it once more.
\end{tcolorbox}
